\begin{theorem}[Artin 代数性公理, 第二版]\label{src:C.7.4}
若下列条件成立, 则 $(\Sch/\Bbbk)_{\text{ét}}$ 上的叠 $\mathcal X$ 为 $\Bbbk$ 上局部有限型的代数叠.
\begin{enumerate}[$(\mathrm{AA}_1)$]
\item (保持极限) 叠 $\mathcal X$ 保持极限.
\item (对角的可表性) 对角 $\mathcal X\to\mathcal X\times\mathcal X$ 可表.
\item (切空间的有限性) 对每个对象 $\xi:\Spec\Bbbk\to\mathcal X$, $\operatorname{Def}_{\xi}(\Bbbk)$ 为有限维 $\Bbbk$-向量空间.
\item (强齐性) 对每个 $\Bbbk$-代数同态 $B_0\to A_0$ 及核的平方为零的 $\Bbbk$-代数满射 $A\twoheadrightarrow A_0$, 函子
\[
\mathcal X(B_0\times_{A_0}A)\to\mathcal X(B_0)\times_{\mathcal X(A_0)}\mathcal X(A)
\]
为等价, 即 $(\mathrm{RS}_4^*)$ 成立.
\item (有效性) 对每个剩余域为 $\Bbbk$ 的完备 Noether 局部 $\Bbbk$-代数 $(R,\mathfrak m)$, 自然函子
\[
\mathcal X(\Spec R)\to\varprojlim\mathcal X(\Spec R/\mathfrak m^n)
\]
为范畴等价.
\item (凝聚形变理论) 对每个 $\Bbbk$-代数 $A$ 上的对象 $\xi\in\mathcal X(A)$, 函子 $\operatorname{Def}_{\xi}(-)$ 与乘积交换.
\item (凝聚障碍理论的存在性) 对每个 $\Bbbk$-代数 $A$ 上的对象 $\xi\in\mathcal X(A)$, 存在如下数据:
\begin{enumerate}[(a)]
\item 存在 $A$-线性函子
\[
\operatorname{Ob}_{\xi}(-):\operatorname{Mod}(A)\to\operatorname{Mod}(A),
\]
且对每个核 $I$ 的平方为零的满射 $B\twoheadrightarrow A$, 存在元素 $\mathrm{ob}_{\xi}(B)\in\operatorname{Ob}_{\xi}(I)$, 使得存在延拓
\[
\begin{tikzcd}
\Spec A\arrow[r,"\xi"]\arrow[d,hook]&\mathcal X\\
\Spec B\arrow[ur,dashed]&
\end{tikzcd}
\]
当且仅当 $\mathrm{ob}_{\xi}(B)=0$, 且
\item 对每个 $\Bbbk$-代数复合 $B\to B'\to A$, 若 $B\twoheadrightarrow A$ 和 $B'\twoheadrightarrow A$ 均满射且核 $I,I'$ 的平方为零, 则 $\mathrm{ob}_{\xi}(B)$ 在 $\operatorname{Ob}_{\xi}(I)\to\operatorname{Ob}_{\xi}(I')$ 下的像为 $\mathrm{ob}_{\xi}(B')$.
\item 对每个 $\Bbbk$-代数 $A$ 上的对象 $\xi\in\mathcal X(A)$, 函子 $\operatorname{Ob}_{\xi}(-)$ 与乘积交换.
\end{enumerate}
\end{enumerate}
此外, $(\mathrm{AA}_2)$ 可替换为
\begin{enumerate}[(1)]
\item[$(\mathrm{AA}_{2'})$] 对每个对象 $\xi:\Spec\Bbbk\to\mathcal X$, $\operatorname{Aut}_{\xi}(\Bbbk)$ 为有限维 $\Bbbk$-向量空间, 且对每个 $\Bbbk$-代数 $A$ 上的对象 $\xi\in\mathcal X(A)$, 函子 $\operatorname{Aut}_{\xi}(-)$ 与乘积交换.
\end{enumerate}
\end{theorem}
\begin{proof}
我们验证定理 \bookref{C.7.1} 的条件. 由 $(\mathrm{AA}_3)$--$(\mathrm{AA}_4)$, 可应用 Rim--Schlessinger 判别法 (\bookref{C.4.6}) 推出万有形式形变的存在性, 即 \bookref{C.7.1}(3) 成立. 剩下只需验证万有性的开性, 即 \bookref{C.7.1}(5). 设 $\xi_0:U_0\to\mathcal X$ 为有限型仿射 $\Bbbk$-概形 $U_0=\Spec B_0$ 出发的态射, 在点 $u_0\in U_0(\Bbbk)$ 处万有. 由 $(\mathrm{AA}_1)$--$(\mathrm{AA}_2)$, 态射 $\xi_0:U_0\to\mathcal X$ 可表且局部有限型. 令 $\Sigma=\{u\in U_0(\Bbbk)\mid\xi_0:U_0\to\mathcal X\text{ 在 }u\text{ 处不万有}\}$. 若万有性的开性不成立, 则 $u_0\in\overline\Sigma$, 且存在由互异点组成的可数无限子集 $\Sigma'=\{u_1,u_2,\ldots\}\subseteq\Sigma$, 满足 $u_0\in\overline{\Sigma'}$.

\textit{步骤 1.} 我们断言存在交换图
\[
\begin{tikzcd}
U_0\arrow[r,hook]\arrow[d,"\xi_0"']&U_1\arrow[r,hook]\arrow[dl,"\xi_1"]&U_2\arrow[r,hook]\arrow[dll,bend left=10,"\xi_2"]&\cdots\\
\mathcal X&&&
\end{tikzcd}
\]
其中每个闭浸入 $U_{n-1}\hookrightarrow U_n$ 由短正合列
\[
0\to\kappa(u_n)\to\mathcal O_{U_n}\to\mathcal O_{U_{n-1}}\to0
\]
定义, 且存在 $u_n$ 的开邻域 $W_n\subseteq U_n\setminus\{u_0,\ldots,u_{n-1}\}$, 使每个限制 $\xi_n|_{W_n}$ 都不是 $\xi_0|_{W_n\cap U_0}$ 的平凡形变. 由构造, 每个 $U_n=\Spec B_n$ 为仿射概形, 每个 $m\geq n$ 时的闭浸入 $U_n\hookrightarrow U_m$ 均由平方为零的理想定义. 假设已构造 $\xi_0,\ldots,\xi_{n-1}$. 因为 $\xi_0:U_0\to\mathcal X$ 和 $\xi_{n-1}:U_{n-1}\to\mathcal X$ 在 $u_n$ 的某个开邻域中同构, 态射 $\xi_{n-1}:U_{n-1}\to\mathcal X$ 在 $u_n$ 处也不万有. 因此存在 $\operatorname{Art}_{\Bbbk}$ 中的满射 $A\to A_0$, 满足 $\ker(A\to A_0)=\Bbbk$, 及交换图
\begin{equation}\label{src:C.7.5}
\begin{tikzcd}
\Spec\Bbbk\arrow[r]\arrow[rr,bend left=25,"u_n"]&\Spec A_0\arrow[r]\arrow[d]&U_{n-1}\arrow[d,"\xi_{n-1}"]\\
&\Spec A\arrow[r]\arrow[ur,dashed,"\nexists"']&\mathcal X,
\end{tikzcd}
\end{equation}
使得 $u_n$ 为 $\Spec A_0\to U_n$ 的像, 且不存在提升 $\Spec A\to U_{n-1}$. 用强齐性 $(\mathrm{AA}_4)$, 可延拓交换图
\[
\begin{tikzcd}[column sep=large]
\Spec A_0\arrow[r]\arrow[d,hook]&U_{n-1}=\Spec B_{n-1}\arrow[d,hook]\arrow[ddr,bend left=20,"\xi_{n-1}"]&\\
\Spec A\arrow[r]\arrow[drr,bend right=15]&U_n=\Spec(A\times_{A_0}B_{n-1})\arrow[dr,dashed,"\xi_n"]&\\
&&\mathcal X,
\end{tikzcd}
\]
得到 $U_n=\Spec B_n$ 上的对象 $\xi_n$, 其中 $B_n:=A\times_{A_0}B_{n-1}$. 若 $\xi_n$ 在 $u_n$ 的开邻域中为 $\xi_0$ 的平凡形变, 则 $\Spec A\to\mathcal X$ 将是 $\Spec A_0$ 的平凡形变, 与 \bookeqref{C.7.5} 不可提升矛盾. 最后注意 $\ker(B_n\to B_{n-1})=\Bbbk$, 因为 $\ker(A\to A_0)=\Bbbk$. 断言得证.

\textit{步骤 2.} 令 $\widehat B=\varprojlim B_n$, $\widehat U=\Spec\widehat B$. 我们断言存在对象 $\widehat\xi\in\mathcal X(\widehat U)$, 延拓每个 $\xi_n\in\mathcal X(U_n)$. 令 $M_n=\ker(B_n\to B_0)$ (注意 $M_0=0$). 因为 $M_n^2=0$, 可将 $M_n$ 看作 $B_0$-模. $\Bbbk$-代数
\[
\widetilde B:=\left\{(b_0,b_1,\ldots)\in\prod_{n\geq0}B_n\ \middle|\ \begin{gathered}\text{每个 }b_n\text{ 在 }B_n\to B_0\\\text{下的像均为 }b_0\end{gathered}\right\}
\]
具有下列性质.
\begin{itemize}
\item 由 $(b_i)\mapsto b_0$ 定义的 $\Bbbk$-代数满同态 $\widetilde B\to B_0$ 的核为 $M:=\prod_{n\geq0}M_n$.
\item 由 $(b_0,b_1,b_2,\ldots)\mapsto(b_0,b_1-b_0,b_2-b_1,b_3-b_2,\ldots)$ 定义的映射 $\widetilde B\to B_0[M]$ 是核的平方为零的 $\Bbbk$-代数满同态.
\item 复合 $\widehat B\to\widetilde B\to B_0[M]$ 诱导短正合列
\[
\begin{tikzcd}[column sep=small]
0\arrow[r]&\ker(\widehat B\to B_0)\arrow[r]\arrow[d,equal]&\ker(\widetilde B\to B_0)\arrow[r]\arrow[d,equal]&\ker(B_0[M]\to B_0)\arrow[r]\arrow[d,equal]&0\\
0\arrow[r]&\varprojlim_{n\geq0}M_n\arrow[r]&\prod_{n\geq0}M_n\arrow[r]&\prod_{n\geq0}M_n\arrow[r]&0.
\end{tikzcd}
\]
其中后一映射为 $(b_0,b_1,\ldots)\mapsto(b_1-b_0,b_2-b_1,\ldots)$.
\item 有等同 $\widehat B=\widetilde B\times_{B_0[M]}B_0$.
\end{itemize}
因为对每个 $n$, $\xi_0$ 的提升 $\xi_n\in\mathcal X(B_n)$ 存在, $\mathrm{ob}_{\xi_0}(B_n)=0\in\operatorname{Ob}_{\xi_0}(M_n)$. 由 $(\mathrm{AA}_7)$(b), 元素 $\mathrm{ob}_{\xi_0}(\widetilde B)$ 在 $\operatorname{Ob}_{\xi_0}(M)\to\operatorname{Ob}_{\xi_0}(M_n)$ 下映到 $\mathrm{ob}_{\xi_0}(B_n)$. 由此, 映射 $\operatorname{Ob}_{\xi_0}(M)\hookrightarrow\prod_n\operatorname{Ob}_{\xi_0}(M_n)$ 单射\footnote{$(\mathrm{AA}_7)$(c) 的假设可减弱为仅要求 $\operatorname{Ob}_{\xi_0}(M)\to\prod_n\operatorname{Ob}_{\xi_0}(M_n)$ 单射, 不过实际中通常验证双射性.}, 因而 $\mathrm{ob}_{\xi_0}(\widetilde B)=0\in\operatorname{Ob}_{\xi_0}(M)$, 表明存在 $\xi_0$ 的提升 $\widetilde\xi\in\mathcal X(\widetilde B)$.

限制 $\widetilde\xi|_{B_n}$ 未必与 $\xi_n$ 同构. 然而, 可用 $\operatorname{Def}_{\xi_0}(M_n)=\operatorname{Lift}_{\xi_0}(B_0[M_n])$ 在非空提升集合 $\operatorname{Lift}_{\xi_0}(\widetilde{B_n})$ 上的自由传递作用, 找到元素 $t_n\in\operatorname{Def}_{\xi_0}(M_n)$, 使 $\xi_n=t_n\cdot\widetilde\xi|_{B_n}$ (引理 \bookref{C.7.3}). 因为由 $(\mathrm{AA}_6)$ 有 $\operatorname{Def}_{\xi_0}(M)\xrightarrow\sim\prod_n\operatorname{Def}_{\xi_0}(M_n)$, 存在映到 $(t_n)$ 的 $\widetilde t\in\operatorname{Def}_{\xi_0}(M)$. 将 $\widetilde\xi$ 替换为 $\widetilde t\cdot\widetilde\xi$ 后, 可使每个 $n$ 的 $\widetilde\xi|_{B_n}$ 与 $\xi_n$ 同构.

现在证明, 经复合 $\widetilde B\to B_0[M]\to B_0[M_n]$ 得到的每个限制 $\widetilde\xi|_{B_0[M_n]}\in\operatorname{Def}_{\xi_0}(M_n)$ 均为平凡形变. 事实上, 映射 $M=\ker(\widetilde B\to B_0)\to\ker(B_0[M_n]\to B_0)=M_n$ 诱导形变模的映射 $\operatorname{Def}_{\xi_0}(M)\to\operatorname{Def}_{\xi_0}(M_n)$, 在 $(\mathrm{AA}_6)$ 的等同 $\operatorname{Def}_{\xi_0}(M)\cong\prod_n\operatorname{Def}_{\xi_0}(M_n)$ 下, 将元素 $(\eta_0,\eta_1,\ldots)$ 映到 $(\eta_{n+1}|_{B_n}-\eta_n)$. 环映射 $\widetilde B\to B_0[M_n]$ 还诱导映射 $\operatorname{Lift}_{\xi_0}(\widetilde B)\to\operatorname{Lift}_{\xi_0}(B_0[M_n])$, 对 $\operatorname{Def}_{\xi_0}(M)\to\operatorname{Def}_{\xi_0}(M_n)$ 等变. 因而 $\widetilde\xi$ 在 $\operatorname{Lift}_{\xi_0}(B_0[M_n])=\operatorname{Def}_{\xi_0}(M_n)$ 中的像为 $\xi_{n+1}|_{B_n}-\xi_n=0$.

现在将等式 $\widehat B=\widetilde B\times_{B_0[M]}B_0$ 与强齐性 $(\mathrm{AA}_4)$ 应用于下图, 即得到延拓 $(\xi_n)\in\varprojlim\mathcal X(B_n)$ 的 $\widehat\xi\in\mathcal X(\widehat B)$ 的存在性:
\[
\begin{tikzcd}[column sep=large]
\Spec B_0[M]\arrow[r]\arrow[d,hook]&\Spec B_0\arrow[d,hook]\arrow[ddr,bend left=20,"\xi_0"]&\\
\Spec\widetilde B\arrow[r]\arrow[drr,bend right=15,"\widetilde\xi"']&\Spec\widehat B\arrow[dr,dashed,"\widehat\xi"]&\\
&&\mathcal X.
\end{tikzcd}
\]
\textit{步骤 3.} 现在用 $\xi_0:U_0\to\mathcal X$ 在 $u_0$ 处的万有性导出矛盾. 因为 $\mathcal X$ 保持极限 $(\mathrm{AA}_1)$, 存在有限生成 $\Bbbk$-子代数 $B'\subseteq\widehat B$, 对象 $\xi'\in\mathcal X(B')$ 及同构 $\widehat\xi\xrightarrow\sim\xi'|_{\widehat B}$. 必要时增大 $B'$, 可假设复合 $B'\hookrightarrow\widehat B\to B_0$ 满射. 因为 $\ker(\widehat B\to B_0)$ 的平方为零, $\ker(B'\to B_0)$ 亦如此. 这定义闭浸入 $U_0\hookrightarrow U':=\Spec B'$, 可考虑交换图
\[
\begin{tikzcd}
U_0\arrow[dr,"i"]\arrow[drr,hook,bend left=15]\arrow[ddr,bend right=25,"\mathrm{id}"']&&\\
&U_0\times_{\mathcal X}U'\arrow[r]\arrow[d]&U'=\Spec B'\arrow[d,"\xi'"]\\
&U_0\arrow[r,"\xi_0"']&\mathcal X,
\end{tikzcd}
\]
其中纤维积 $U_0\times_{\mathcal X}U'$ 是 $\Bbbk$ 上局部有限型的代数空间. 因为 $\xi_0:U_0\to\mathcal X$ 在 $u_0$ 处万有, 由光滑性的无穷小提升判别法的 Artin 版本 (\bookref{4.7.1}), $U_0\times_{\mathcal X}U'\to U'$ 在 $i(u_0)$ 处光滑. 将 $U_0,U'$ 替换为仿射开邻域, 并将 $\{u_1,u_2,\ldots\}$ 替换为包含于这些开子集中的无限子序列后, 可使 $U_0\times_{\mathcal X}U'\to U'$ 光滑. 概形的光滑性的无穷小提升判别法的非 Artin 版本 (\bookref{A.3.1}) 蕴含, $U_0\times_{\mathcal X}U'\to U'$ 在 $U_0$ 上的截面延拓为整体截面 $U'\to U_0\times_{\mathcal X}U'$. 这蕴含 $\xi'$ 为 $\xi_0$ 的平凡形变, 进而每个 $\xi_n$ 都是 $\xi_0$ 的平凡形变, 矛盾.

对于补充断言, 我们证明 $(\mathrm{AA}_{2'})$ 蕴含 $(\mathrm{AA}_2)$, 即对角的可表性. 用 \S\bookref{B.3} 的极限方法, 只需考虑有限型 $\Bbbk$-概形出发的映射 $(a,b):T\to\mathcal X\times_{\Bbbk}\mathcal X$. 基变换 $\underline{\Isom}_T(a,b)\to T$ 的形变函子 $\operatorname{Def}_{\xi}(-)$ 对应于 $\mathcal X$ 的自同构函子 $\operatorname{Aut}_{\xi}(-)$, 并取 $\mathcal X$ 的形变函子 $\operatorname{Def}_{\xi}(-)$ 作为 $\underline{\Isom}_T(a,b)$ 的障碍函子 $\operatorname{Ob}_{\xi}(-)$. 我们的叙述遵循 \cite[Tag 0CYF]{SP} 和 \cite[定理 A]{Hal17}. 另见 \cite{Art69b}, \cite[定理 $5.3$]{Art74b} 及 \cite[主定理]{HR19a} 中 Artin 判别法的其他版本, 以及 \cite[定理 $1$]{Mur95} 和 \cite[定理 $1$]{Mur64} 中取值于阿贝尔群的函子的可表性判别法.
\end{proof}
\begin{remark}\label{src:C.7.6}
本定理的逆命题也成立. 为证明条件的必要性, 只需验证 $(\mathrm{AA}_3)$, $(\mathrm{AA}_4)$, $(\mathrm{AA}_6)$ 和 $(\mathrm{AA}_7)$. 条件 $(\mathrm{AA}_3)$ (切空间的有限性) 成立, 因为 $\mathcal X$ 在 $\Bbbk$ 上有限型; $(\mathrm{AA}_4)$ (强齐性) 由 \cite[Tag 07WN]{SP} 成立. 条件 $(\mathrm{AA}_7)$ (障碍理论的存在性) 来自 $\mathcal X$ 的余切复形 $\mathrm L_{\mathcal X/\Bbbk}$ 的存在性, 它满足类似定理 \bookref{C.3.1} 的性质; 见 \cite{Ols06}. 若 $\xi:\Spec A\to\mathcal X$ 为有限生成 $\Bbbk$-代数 $A$ 出发的态射, $I$ 为 $A$-模, 则令 $\operatorname{Ob}_{\xi}(I):=\Ext_A^1(\xi^*L_{\mathcal X/\Bbbk},I)$. 最后, 由上同调与基变换 (见 \cite{Hal14a}), $\operatorname{Def}_{\xi}(-)$ 和 $\operatorname{Ob}_{\xi}(-)$ 与乘积交换.
\end{remark}
